\chapter{The Bootstrap}
\label{ch:bootstrap}

Chapter 7 built confidence intervals for a mean and a proportion. Each one
rested on a formula for the standard error: $s/\sqrt{n}$ for a mean,
$\sqrt{\hat p(1-\hat p)/n}$ for a proportion. Those formulas came from Chapter 6,
where we worked out by algebra how much the sample mean wobbles from sample to
sample.

Now your dashboard shows a median latency of 41 ms, a p90 of 71 ms, and a
week-over-week drop in the median of 17 ms. Your manager asks: ``plus or minus
what?'' There is no $s/\sqrt{n}$ for a median. There is no tidy formula for
the standard error of a p90, of a ratio of two means, or of a correlation
coefficient under a skewed population. The algebra that worked for the mean
either gets ugly or needs things you do not know, like the height of the
population density at the median.

The bootstrap replaces that algebra with computation. It answers ``how much
would this number wobble if I could re-run the sampling?'' by re-running the
sampling on the data itself. This chapter shows how it works, when to trust it,
when an exact method beats it, and where it breaks.

\section{The problem: statistics with no formula}
\label{sec:boot-problem}

Here are 40 request latencies (ms) from one service, sorted:
\begin{center}
\small
\begin{tabular}{*{10}{r}}
12 & 18 & 21 & 23 & 24 & 24 & 24 & 24 & 25 & 27\\
29 & 31 & 31 & 32 & 33 & 36 & 36 & 37 & 39 & 40\\
42 & 43 & 48 & 50 & 50 & 51 & 52 & 53 & 55 & 58\\
59 & 63 & 64 & 68 & 68 & 70 & 79 & 79 & 83 & 145
\end{tabular}
\end{center}
The mean is $46.15$, the median is $(40+42)/2 = 41$, and the sample SD is
$24.48$. Figure~\ref{fig:bootstrap-dotplot} shows the shape: a hump in the 20s
to 50s and a long right tail ending at one 145 ms request.

\begin{figure}[H]
\centering
\begin{tikzpicture}
\begin{axis}[
  width=12cm, height=4.6cm,
  xmin=0, xmax=155, ymin=0, ymax=7.5,
  axis lines=left, ytick={0,2,4,6}, xtick={0,20,...,140},
  xlabel={latency (ms), grouped in 5 ms bins}, ylabel={count},
  clip=false,
]
\addplot[only marks, mark=*, mark size=2.6pt, color=histfill, draw=distblue]
coordinates {(12.5,1) (17.5,1) (22.5,1) (22.5,2) (22.5,3) (22.5,4) (22.5,5) (22.5,6)
(27.5,1) (27.5,2) (27.5,3) (32.5,1) (32.5,2) (32.5,3) (32.5,4) (37.5,1) (37.5,2)
(37.5,3) (37.5,4) (42.5,1) (42.5,2) (42.5,3) (47.5,1) (52.5,1) (52.5,2) (52.5,3)
(52.5,4) (52.5,5) (57.5,1) (57.5,2) (57.5,3) (62.5,1) (62.5,2) (67.5,1) (67.5,2)
(72.5,1) (77.5,1) (77.5,2) (82.5,1) (147.5,1)};
\draw[densityred, thick] (axis cs:41,0) -- (axis cs:41,6.6);
\node[densityred, anchor=south, font=\small] at (axis cs:41,6.6) {median 41};
\draw[targetgreen, thick, dashed] (axis cs:46.15,0) -- (axis cs:46.15,5.6);
\node[targetgreen, anchor=west, font=\small] at (axis cs:47,5.9) {mean 46.15};
\end{axis}
\end{tikzpicture}
\caption{Forty latencies. The right tail pulls the mean above the median. We
want an interval for the population median, and later for the p90, and no
textbook formula gives one directly.}
\label{fig:bootstrap-dotplot}
\end{figure}

For the mean we already know the answer: $\SE = 24.48/\sqrt{40} = 3.87$, and
Chapter 7 gives a $t$-interval. For the median there is a large-sample formula,
\[
  \SE(\text{median}) \approx \frac{1}{2 f(m)\sqrt{n}},
\]
where $f(m)$ is the height of the population density at the true median $m$.
A \emph{density} is the smooth curve whose area over an interval gives the
probability of landing there (Chapter 1). We do not know $f$, and estimating its
height from 40 points is shaky. For a p90, a ratio of two means, or a
correlation, the formulas get worse. We need a method that works for
\emph{any} statistic you can compute.

\begin{recoverybox}
Pause and think: what single object would let you compute the standard error
of \emph{any} statistic by brute force? (Answer: the population itself. If you
had it, you could draw thousands of fresh samples and watch the statistic vary.
The bootstrap's trick is a stand-in for that population.)
\end{recoverybox}

\section{The plug-in principle}
\label{sec:boot-plugin}

Every standard error is a property of the \emph{population} $F$ (the full
distribution of latencies you would see if you logged forever) and the sample
size $n$. If we knew $F$, we would simulate: draw $n$ values from $F$, compute
the median, repeat 10{,}000 times, and take the SD of the 10{,}000 medians.
That SD \emph{is} the standard error. This is the sampling distribution of
Chapter 6, run by computer.

We do not know $F$. The plug-in principle says: use your best estimate of $F$
in its place, and run the same simulation.

\begin{definition}[Empirical distribution]
The \textbf{empirical distribution} $\hat F_n$ of data $x_1,\dots,x_n$ puts
probability $1/n$ on each data point. Picture: a bag containing your $n$
values written on slips of paper. Its cumulative distribution function, the
\textbf{ECDF}, is
\[
  \hat F_n(t) = \frac{\#\{i : x_i \le t\}}{n},
\]
the fraction of data at or below $t$. It is a staircase that rises by $1/n$ at
each data point.
\end{definition}

\begin{figure}[H]
\centering
\begin{tikzpicture}
\begin{axis}[
  width=12cm, height=6cm,
  xmin=0, xmax=160, ymin=0, ymax=1.05,
  axis lines=left, xlabel={latency $t$ (ms)}, ylabel={fraction $\le t$},
  legend pos=south east, legend style={font=\small},
]
\addplot[densityred, thick, smooth] coordinates {(5,0.000) (10,0.010) (15,0.051)
(20,0.124) (25,0.217) (30,0.316) (35,0.412) (40,0.500) (45,0.578) (50,0.645)
(55,0.702) (60,0.750) (65,0.791) (70,0.825) (75,0.853) (80,0.876) (85,0.895)
(90,0.912) (95,0.925) (100,0.937) (105,0.946) (110,0.954) (115,0.961) (120,0.966)
(125,0.971) (130,0.975) (135,0.979) (140,0.982) (145,0.984) (150,0.986)
(155,0.988) (160,0.990)};
\addlegendentry{true population CDF $F$}
\addplot[distblue, thick, const plot] coordinates {(0,0) (12,0.025) (18,0.050)
(21,0.075) (23,0.100) (24,0.200) (25,0.225) (27,0.250) (29,0.275) (31,0.325)
(32,0.350) (33,0.375) (36,0.425) (37,0.450) (39,0.475) (40,0.500) (42,0.525)
(43,0.550) (48,0.575) (50,0.625) (51,0.650) (52,0.675) (53,0.700) (55,0.725)
(58,0.750) (59,0.775) (63,0.800) (64,0.825) (68,0.875) (70,0.900) (79,0.950)
(83,0.975) (145,1.000) (160,1.000)};
\addlegendentry{ECDF $\hat F_{40}$ of our sample}
\end{axis}
\end{tikzpicture}
\caption{The staircase ECDF of the 40 latencies hugs the true CDF (these data
were simulated from a known log-normal population so we can see both). The
bootstrap uses the staircase as if it were the population.}
\label{fig:bootstrap-ecdf}
\end{figure}

Figure~\ref{fig:bootstrap-ecdf} shows why this is reasonable. The law of large
numbers (Probability From the Ground Up, Ch.~20) says each staircase height
$\hat F_n(t)$ converges to $F(t)$, and a stronger result (the
Glivenko--Cantelli theorem) says the whole staircase converges uniformly. With
40 points the staircase is already close.

\begin{intuitionbox}
The plug-in principle is something you already do. You estimate the
population mean by the sample mean, which is the mean of $\hat F_n$. You
estimate the population median by the sample median, the median of $\hat F_n$.
The bootstrap goes one step further: it estimates the \emph{sampling
variability} of a statistic under $F$ by its sampling variability under
$\hat F_n$.
\end{intuitionbox}

Figure~\ref{fig:bootstrap-worlds} lays out the parallel.

\begin{figure}[H]
\centering
\begin{tikzpicture}[
  box/.style={draw, rounded corners, minimum width=2.6cm, minimum height=0.9cm,
    align=center, font=\small},
  arr/.style={-{Stealth[length=2.5mm]}, thick},
]
\node[font=\bfseries\small, sourceblue] at (0,1.2) {Real world};
\node[box, fill=sourceblue!10] (F) at (0,0) {population $F$\\(unknown)};
\node[box, fill=sourceblue!10] (x) at (0,-1.7) {sample\\$x_1,\dots,x_n$};
\node[box, fill=sourceblue!10] (t) at (0,-3.4) {statistic\\$\hat\theta = s(x)$};
\draw[arr] (F) -- node[right, font=\footnotesize] {draw $n$} (x);
\draw[arr] (x) -- node[right, font=\footnotesize] {compute} (t);
\node[font=\bfseries\small, bootorange] at (6.5,1.2) {Bootstrap world};
\node[box, fill=bootorange!15] (Fh) at (6.5,0) {empirical $\hat F_n$\\(known)};
\node[box, fill=bootorange!15] (xs) at (6.5,-1.7) {resample\\$x^*_1,\dots,x^*_n$};
\node[box, fill=bootorange!15] (ts) at (6.5,-3.4) {replicate\\$\hat\theta^* = s(x^*)$};
\draw[arr] (Fh) -- node[right, font=\footnotesize] {draw $n$ with replacement} (xs);
\draw[arr] (xs) -- node[right, font=\footnotesize] {compute} (ts);
\draw[arr, dashed, gray] (x) -- node[above, font=\footnotesize] {becomes} (Fh);
\node[font=\footnotesize, align=center, text width=4.2cm] at (0,-4.6)
  {spread of $\hat\theta$ around $\theta$: what we want, cannot see};
\node[font=\footnotesize, align=center, text width=4.2cm] at (6.5,-4.6)
  {spread of $\hat\theta^*$ around $\hat\theta$: repeat $B$ times and look};
\end{tikzpicture}
\caption{The bootstrap copies the real-world sampling process inside a world
we fully control. The sample plays the role of the population; resamples play
the role of samples.}
\label{fig:bootstrap-worlds}
\end{figure}

\section{Resampling with replacement, by hand}
\label{sec:boot-resample}

Drawing from $\hat F_n$ means: pick one of the $n$ data points uniformly at
random, write it down, put it back, and repeat $n$ times. The result is a
\textbf{bootstrap resample}: a list of $n$ values, each one copied from the
data, where some originals appear twice or more and some not at all.

Take a tiny dataset of $n=8$ latencies:
\[
  12,\ 15,\ 18,\ 21,\ 24,\ 30,\ 45,\ 110 .
\]
The sample median is $(21+24)/2 = 22.5$. Figure~\ref{fig:bootstrap-resamples}
shows five random resamples, each sorted so the median (average of positions 4
and 5) is easy to read.

\begin{figure}[H]
\centering
\begin{tikzpicture}[
  cell/.style={draw, minimum width=0.82cm, minimum height=0.55cm, font=\small,
    inner sep=0pt},
]
\def\rowlabel#1#2{\node[anchor=east, font=\small] at (-0.2,#1) {#2};}
\rowlabel{0}{data}
\foreach \v [count=\i] in {12,15,18,21,24,30,45,110} {
  \ifnum\v=110 \node[cell, fill=trapred!20] at (0.82*\i,0) {\v};
  \else \node[cell, fill=sourceblue!12] at (0.82*\i,0) {\v}; \fi }
\node[anchor=west, font=\small] at (7.8,0) {median $= 22.5$};
\draw[gray] (0.2,-0.5) -- (10.6,-0.5);
\foreach \row/\vals/\med [count=\r] in {
  1/{12,15,21,24,45,45,110,110}/34.5,
  2/{15,18,18,18,21,21,45,45}/19.5,
  3/{12,12,24,24,30,45,45,45}/27,
  4/{12,12,18,18,21,21,45,45}/19.5,
  5/{15,15,18,21,21,24,110,110}/21} {
  \rowlabel{-0.75*\r-0.3}{resample \row}
  \foreach \v [count=\i] in \vals {
    \ifnum\v=110 \node[cell, fill=trapred!20] at (0.82*\i,-0.75*\r-0.3) {\v};
    \else \node[cell, fill=bootorange!15] at (0.82*\i,-0.75*\r-0.3) {\v}; \fi }
  \draw[densityred, thick] (0.82*4-0.41,-0.75*\r-0.3-0.275) rectangle
    (0.82*5+0.41,-0.75*\r-0.3+0.275);
  \node[anchor=west, font=\small] at (7.8,-0.75*\r-0.3) {median $= \med$};
}
\end{tikzpicture}
\caption{Five bootstrap resamples of the 8-point dataset (sorted; the red frame
marks the two middle values). The outlier 110 appears twice in resamples 1 and
5 and not at all in 2, 3, 4. The resample medians, $34.5, 19.5, 27, 19.5, 21$,
spread around the original $22.5$, and that spread is what the bootstrap
measures.}
\label{fig:bootstrap-resamples}
\end{figure}

Three facts are visible already.
\begin{itemize}
  \item Each resample has exactly $n=8$ values. Same size as the data, because
  we are imitating a sample of size $n$.
  \item Duplicates are the whole point. Sampling \emph{without} replacement
  would give back the same 8 numbers in a new order and the same median every
  time.
  \item The chance that a particular data point is missing from a resample is
  $(1-1/n)^n$. For $n=8$ that is $(7/8)^8 = 0.344$; as $n$ grows it tends to
  $1/e \approx 0.368$. So a typical resample uses about $63\%$ of the distinct
  original points.
\end{itemize}

The five medians have SD $6.49$. With five resamples that is a crude
estimate, but it is already a standard error for the median of 8 latencies,
computed with no formula. Run $B = 2000$ resamples instead of 5 and the
estimate stabilises.

\begin{trapbox}
Resample the \emph{units} that were sampled independently. For a correlation
between $x$ and $y$, each resample draws whole $(x_i, y_i)$ pairs. Resampling
the $x$'s and the $y$'s separately destroys the pairing: every replicate is a
correlation between unrelated numbers, so the replicates cluster around zero
whatever the real correlation is.
\end{trapbox}

\section{The bootstrap distribution}
\label{sec:boot-distribution}

Back to the 40 latencies. Draw $B = 2000$ resamples of size 40, compute the
median of each, and plot the 2000 medians.

\begin{lstlisting}[language=Python]
import numpy as np
lat = np.array([48,21,63,70,12,18,43,33,40,24,68,64,42,79,53,24,50,23,68,39,
                36,27,83,36,31,32,55,50,51,52,145,31,29,25,58,79,37,24,24,59])
rng = np.random.default_rng(0)
boot = rng.choice(lat, size=(2000, 40), replace=True)  # 2000 resamples
med_star = np.median(boot, axis=1)                      # 2000 medians
print(med_star.std(ddof=1))                   # bootstrap SE   -> 5.28
print(np.quantile(med_star, [0.025, 0.975]))  # percentile CI  -> [33. 52.]
\end{lstlisting}

\begin{figure}[H]
\centering
\begin{tikzpicture}
\begin{axis}[
  width=12cm, height=6.2cm,
  xmin=27, xmax=63, ymin=0, ymax=400,
  axis lines=left, xlabel={bootstrap median $\hat\theta^*$ (ms)},
  ylabel={count (of 2000)}, xtick={29,33,...,61},
  clip=false,
]
\addplot[ybar interval, fill=bootorange!60, draw=bootorange!80!black]
coordinates {(27,0) (29,3) (31,46) (33,67) (35,237) (37,270) (39,333) (41,338)
(43,108) (45,170) (47,74) (49,261) (51,76) (53,12) (55,4) (57,0) (59,1) (61,0)};
\draw[densityred, very thick] (axis cs:41,0) -- (axis cs:41,370);
\node[densityred, anchor=south, font=\small] at (axis cs:41,370) {$\hat\theta = 41$};
\draw[distblue, thick, dashed] (axis cs:33,0) -- (axis cs:33,300);
\draw[distblue, thick, dashed] (axis cs:52,0) -- (axis cs:52,300);
\draw[distblue, thick, {Stealth}-{Stealth}] (axis cs:33,300) -- (axis cs:52,300);
\node[distblue, anchor=south west, font=\small] at (axis cs:52.3,300)
  {95\% percentile interval};
\node[distblue, anchor=north, font=\small] at (axis cs:33,-30) {33};
\node[distblue, anchor=north, font=\small] at (axis cs:52,-30) {52};
\end{axis}
\end{tikzpicture}
\caption{Bootstrap distribution of the median from 2000 resamples (bins 2 ms
wide). The middle $95\%$ of the bars runs from 33 to 52 ms. The comb-like
lumps are real: a median of 40 values can only be an average of two data
values, so it takes few distinct values.}
\label{fig:bootstrap-median-hist}
\end{figure}

\begin{definition}[Bootstrap distribution]
The \textbf{bootstrap distribution} of a statistic is the distribution of
$\hat\theta^* = s(x^*)$ when $x^*$ is a resample from the data. In practice we
see it through $B$ replicates $\hat\theta^*_1, \dots, \hat\theta^*_B$. It is our
estimate of the sampling distribution of $\hat\theta$.
\end{definition}

The bootstrap distribution estimates the \emph{shape and spread} of the
sampling distribution. It does not estimate where the truth is. Its centre sits
near $\hat\theta$, not near $\theta$: the mean of the 2000 medians is $41.84$,
close to the sample median 41. Nothing in resampling can move you toward a
population value you never observed.

\begin{happeningbox}
The bootstrap distribution measures how much the statistic changes when the
sample is perturbed in the way random sampling perturbs it. Picture shaking the
data: some points double up, some vanish. A statistic that barely moves under
this shaking (the median of 10{,}000 values) has a small SE. A statistic that
jumps around (the max of 40 values) has a large or meaningless one.
\end{happeningbox}

\section{Bootstrap standard error}
\label{sec:boot-se}

\begin{definition}[Bootstrap standard error]
\[
  \widehat{\SE}_{\text{boot}} = \sqrt{\frac{1}{B-1}\sum_{b=1}^{B}
  \left(\hat\theta^*_b - \bar\theta^*\right)^2},
  \qquad \bar\theta^* = \frac{1}{B}\sum_{b=1}^B \hat\theta^*_b .
\]
It is the ordinary sample SD (Chapter 2) of the $B$ replicates.
\end{definition}

For the latency median, $\widehat{\SE}_{\text{boot}} = 5.28$ ms. The mean's
formula SE is $3.87$ ms. Is the median the noisier statistic here? Because
these data were simulated, we know the population and can find the true SEs by
brute force (draw 200{,}000 fresh samples of 40): $4.7$ ms for the median and
$5.0$ ms for the mean. Both one-sample estimates are off. The bootstrap
overstates the median's SE because the values near the middle of this sample
happen to be spread out ($39, 40, 42, 43, 48$), so reshuffles move the midpoint
by several ms. The formula understates the mean's SE because 40 draws rarely
include the population's rare very slow requests, so $s$ comes out too small.
An estimated SE is itself a noisy number computed from one sample.

\begin{whybox}
Why should the SD of replicates estimate the SE? The SE is a fixed function
of the population and $n$: write it $\SE(F, n)$. The bootstrap SE is the same
function evaluated at the empirical distribution, $\SE(\hat F_n, n)$, which is
the plug-in principle again. As $\hat F_n \to F$, the two agree for any
statistic that depends smoothly on $F$. The $B$ replicates are only a Monte
Carlo approximation to $\SE(\hat F_n, n)$, and Section~\ref{sec:boot-B} shows how
big $B$ must be.
\end{whybox}

A quick sanity check uses the mean. One draw from $\hat F_n$ is a random data
point, so its variance is $\hat\sigma^2 = \frac1n\sum_i (x_i - \xbar)^2$, the
SD computed with divisor $n$ instead of $n-1$. A resample mean averages $n$
independent such draws, so by Chapter 6 its SD is $\hat\sigma/\sqrt n$. With
unlimited resamples the bootstrap SE of the mean is exactly this:
$24.17/\sqrt{40} = 3.82$, against the formula's $3.87$. (Our 2000 resamples
give $3.86$; the gap from $3.82$ is Monte Carlo noise,
Section~\ref{sec:boot-B}.)

A normal-theory interval $\hat\theta \pm 1.96\,\widehat{\SE}_{\text{boot}}$
is the simplest use: $41 \pm 1.96(5.28) = [30.7, 51.3]$. It assumes the
sampling distribution is bell-shaped and centred at $\theta$. The bootstrap can
do better by using the shape it sees.

\section{The percentile interval}
\label{sec:boot-percentile}

\begin{definition}[Percentile interval]
The $95\%$ \textbf{bootstrap percentile interval} is
\[
  \left[\, \hat\theta^*_{(0.025)},\ \hat\theta^*_{(0.975)} \,\right],
\]
the $2.5$th and $97.5$th percentiles of the $B$ replicates. Picture: sort the
2000 replicates and chop off 50 from each end.
\end{definition}

For the median latency this is $[33, 52]$ ms. The interval is not symmetric
around 41: it reaches 8 ms below and 11 ms above, because the bootstrap
distribution has more room on the right, which mirrors the right-skewed data.

\begin{whybox}
Why should chopping the tails of the bootstrap histogram give a confidence
interval? Suppose some increasing transformation $g$ (a log, a square root)
makes things symmetric: $g(\hat\theta)$ is normal around $g(\theta)$ with some
SD $\tau$, and in the bootstrap world $g(\hat\theta^*)$ is normal around
$g(\hat\theta)$ with the same $\tau$. Follow three steps.
\begin{enumerate}
  \item On the $g$ scale the ordinary interval is
  $g(\hat\theta) \pm 1.96\tau$, and it covers $g(\theta)$ $95\%$ of the time.
  \item The bootstrap replicates $g(\hat\theta^*)$ are normal with centre
  $g(\hat\theta)$ and SD $\tau$, so their $2.5$th and $97.5$th percentiles are
  exactly $g(\hat\theta) - 1.96\tau$ and $g(\hat\theta) + 1.96\tau$.
  \item Percentiles pass through increasing functions unchanged: the
  $2.5$th percentile of $g(\hat\theta^*)$ is $g$ of the $2.5$th percentile of
  $\hat\theta^*$. Mapping the interval back with $g^{-1}$ therefore gives
  $[\hat\theta^*_{(0.025)}, \hat\theta^*_{(0.975)}]$.
\end{enumerate}
So the percentile interval is the right interval whenever such a $g$ exists,
and you never need to know $g$. That is why it handles skewness reasonably
well and why it respects natural ranges: the interval for a proportion never
leaves $[0,1]$.
\end{whybox}

\section{The basic interval}
\label{sec:boot-basic}

There is a second way to read the same histogram. Suppose the error
$\hat\theta - \theta$ in the real world behaves like the error
$\hat\theta^* - \hat\theta$ in the bootstrap world. The bootstrap errors run
from $33 - 41 = -8$ to $52 - 41 = +11$ ($95\%$ of the time). So
\[
  -8 \le \hat\theta - \theta \le 11
  \quad\Longleftrightarrow\quad
  \hat\theta - 11 \le \theta \le \hat\theta + 8 .
\]

\begin{definition}[Basic interval]
\[
  \left[\, 2\hat\theta - \hat\theta^*_{(0.975)},\ 2\hat\theta - \hat\theta^*_{(0.025)} \,\right].
\]
\end{definition}

For the median latency: $[82 - 52,\ 82 - 33] = [30, 49]$. The basic interval
is the percentile interval reflected through $\hat\theta$
(Figure~\ref{fig:bootstrap-pct-basic}).

\begin{figure}[H]
\centering
\begin{tikzpicture}[x=0.32cm]
\draw[-{Stealth}] (25,0) -- (60,0) node[right, font=\small] {ms};
\foreach \t in {25,30,...,55} { \draw (\t,0.1) -- (\t,-0.1)
  node[below, font=\footnotesize] {\t}; }
\draw[densityred, very thick] (41,-0.3) -- (41,2.4);
\node[densityred, above, font=\small] at (41,2.4) {$\hat\theta=41$};
\draw[distblue, line width=2.5pt] (33,1.6) -- (52,1.6);
\node[distblue, anchor=west, font=\small] at (52.5,1.6) {percentile $[33,52]$};
\draw[quotientorange, line width=2.5pt] (30,0.8) -- (49,0.8);
\node[quotientorange, anchor=east, font=\small] at (29.5,0.8) {basic $[30,49]$};
\node[distblue, above, font=\footnotesize] at (37,1.65) {8};
\node[distblue, above, font=\footnotesize] at (46.5,1.65) {11};
\node[quotientorange, below, font=\footnotesize] at (35.5,0.75) {11};
\node[quotientorange, below, font=\footnotesize] at (45,0.75) {8};
\end{tikzpicture}
\caption{The basic interval is the percentile interval flipped around
$\hat\theta$: the long arm on the right becomes a long arm on the left. Which
one is right depends on whether the skew in the bootstrap distribution is a
property of the statistic's sampling distribution (percentile) or an error to
be corrected (basic).}
\label{fig:bootstrap-pct-basic}
\end{figure}

The two readings disagree whenever the bootstrap distribution is skewed. The
basic interval is right if the bootstrap world reproduces the real-world
\emph{error} faithfully, shape and all. The percentile interval is right if a
hidden symmetrising transformation exists. For the median of a skewed
distribution, the simulation in Section~\ref{sec:boot-coverage} shows the
percentile interval covering much better. In practice the basic interval is
seldom the best choice; it is worth knowing because it is the logic behind the
next two methods.

\section{Better intervals: bootstrap-t and BCa}
\label{sec:boot-t-bca}

Percentile and basic intervals are \textbf{first-order accurate}: the chance
that the truth falls beyond one particular endpoint misses its target $2.5\%$
by an error that shrinks like $1/\sqrt{n}$. Two refinements shrink that error
to order $1/n$ for smooth statistics like means, ratios and correlations.

\subsection{Bootstrap-t}

The $t$-interval of Chapter 7 works because $(\bar X - \mu)/\SE$ has a
distribution (Student's $t$) that does not depend on the unknowns. A quantity
like that, whose distribution is free of unknown parameters, is called a
\textbf{pivot}. The bootstrap-t builds a pivot for any statistic and lets the
bootstrap find its distribution instead of looking it up in a $t$ table:
\[
  t^*_b = \frac{\hat\theta^*_b - \hat\theta}{\widehat{\SE}^*_b}, \qquad
  \text{interval} = \left[\, \hat\theta - t^*_{(0.975)}\,\widehat{\SE},\
  \hat\theta - t^*_{(0.025)}\,\widehat{\SE} \,\right].
\]
Here $\widehat{\SE}^*_b$ is the standard error computed \emph{inside}
resample $b$. For a mean that is just $s^*_b/\sqrt n$. For a statistic with no
formula, it requires a second, nested bootstrap inside each resample: $B=1000$
outer times $200$ inner means $200{,}000$ medians. Studentizing (dividing by an
estimated SE) removes the dependence of spread on $\theta$, which is where the
extra accuracy comes from.

Prefer bootstrap-t when the statistic has a stable SE estimate: means, trimmed
means, differences of means, regression slopes. Avoid it for the median and
other quantiles, whose resample SEs are themselves erratic; the ratio then
explodes when $\widehat{\SE}^*_b$ is tiny.

\subsection{BCa: bias-corrected and accelerated}

BCa keeps the percentile idea but picks different percentiles. Two
adjustments are estimated from the data.
\begin{itemize}
  \item \textbf{Bias correction} $z_0$: if more than half the replicates fall
  below $\hat\theta$, the bootstrap distribution is shifted relative to
  $\hat\theta$. Then $z_0 = \Phi^{-1}(\text{fraction of } \hat\theta^*_b < \hat\theta)$,
  where $\Phi^{-1}$ turns a fraction into a standard normal quantile.
  \item \textbf{Acceleration} $a$: measures how fast the SE changes as
  $\theta$ changes, estimated from the jackknife (recompute the statistic
  leaving out one point at a time; $a$ is proportional to the skewness of
  those $n$ values).
\end{itemize}
The interval uses replicate percentiles $\alpha_1, \alpha_2$ instead of
$0.025, 0.975$:
\[
  \alpha_{1,2} = \Phi\!\left( z_0 + \frac{z_0 + z_{\alpha}}{1 - a\,(z_0 + z_{\alpha})} \right),
  \qquad z_\alpha = \mp 1.96 .
\]
With $z_0 = a = 0$ this collapses to the percentile interval.

\begin{lstlisting}[language=Python]
from scipy import stats
res = stats.bootstrap((lat,), np.median, n_resamples=2000,
                      method='BCa', random_state=0)
print(res.confidence_interval)      # -> (32.7, 51.5)
p90 = lambda a, axis: np.quantile(a, 0.9, axis=axis)
res = stats.bootstrap((lat,), p90, n_resamples=2000,
                      method='BCa', random_state=0)
print(res.confidence_interval)      # -> (63.5, 89.2)
\end{lstlisting}

\begin{trapbox}
BCa's acceleration comes from the jackknife, and the jackknife is unreliable
for the median: leaving out one of 40 points moves the median to one of only
two values. For the median, BCa is close to the percentile interval and gains
little. Its real advantage is for smooth statistics (means of skewed data,
ratios, correlations, variances).
\end{trapbox}

\section{Worked examples: median and p90 latency}
\label{sec:boot-worked}

\begin{walkbox}
\textbf{Question.} From the 40 latencies, give a $95\%$ interval for the
population median and for the population p90.

\textbf{Median.} $\hat\theta = 41$. With $B=2000$ (seed 0):
$\widehat{\SE}_{\text{boot}} = 5.28$; percentile $[33, 52]$; basic $[30, 49]$;
BCa $[32.7, 51.5]$. Report: ``median 41 ms, $95\%$ CI 33 to 52 ms.''

\textbf{p90.} The sample p90 (\texttt{numpy}'s default linear interpolation,
Chapter 3) is $70.9$. From the same 2000 resamples:
$\widehat{\SE}_{\text{boot}} = 8.35$; percentile $[59.4, 83.0]$;
basic $[141.8 - 83.0,\ 141.8 - 59.4] = [58.8, 82.4]$; BCa $[63.5, 89.2]$.

The p90 intervals disagree more than the median intervals, and the BCa interval
is shifted right. That disagreement is a warning sign: the tail of 40 points is
thin, and Figure~\ref{fig:bootstrap-p90-hist} shows why.
\end{walkbox}

\begin{figure}[H]
\centering
\begin{tikzpicture}
\begin{axis}[
  width=12cm, height=6cm,
  xmin=50, xmax=150, ymin=0, ymax=700,
  axis lines=left, xlabel={bootstrap p90 $\hat\theta^*$ (ms)},
  ylabel={count (of 2000)}, xtick={50,60,...,150},
  clip=false,
]
\addplot[ybar interval, fill=bootorange!60, draw=bootorange!80!black]
coordinates {(50,5) (55,82) (60,221) (65,561) (70,424) (75,597) (80,71) (85,29)
(90,0) (95,0) (100,0) (105,0) (110,0) (115,0) (120,0) (125,0) (130,0) (135,0)
(140,0) (145,10) (150,0)};
\draw[densityred, very thick] (axis cs:70.9,0) -- (axis cs:70.9,650);
\node[densityred, anchor=south, font=\small] at (axis cs:70.9,650) {$\hat\theta=70.9$};
\draw[distblue, thick, {Stealth}-{Stealth}] (axis cs:59.4,450) -- (axis cs:83,450);
\node[distblue, anchor=west, font=\small] at (axis cs:84,450) {percentile $[59.4, 83.0]$};
\node[anchor=south, font=\small, align=center] at (axis cs:140,40)
  {10 resamples drew the\\145 ms request $\ge 5$ times};
\end{axis}
\end{tikzpicture}
\caption{Bootstrap distribution of the p90. Almost all replicates are
interpolations between the handful of data values above 63 ms, so the
histogram is lumpy and its right edge is pinned to the largest observations.
The bootstrap cannot imagine a latency above 145 ms.}
\label{fig:bootstrap-p90-hist}
\end{figure}

The p90 replicates are built from the top 5 or so data points. A resample can
only rearrange those points; it cannot produce the 100 ms request that the
population will occasionally send. The further into the tail you go, the less
the bootstrap knows. For p99 from 40 points it knows essentially nothing.

\section{Difference in medians between two groups}
\label{sec:boot-two-groups}

After a deploy, 35 latencies from the new version have median 24 ms, against
41 ms for the 40 old ones (Figure~\ref{fig:bootstrap-two-box}).

\begin{figure}[H]
\centering
\begin{tikzpicture}
\begin{axis}[
  width=12cm, height=4.4cm,
  xmin=0, xmax=155, ymin=0.4, ymax=2.6,
  ytick={1,2}, yticklabels={new ($n=35$), old ($n=40$)},
  xlabel={latency (ms)}, axis x line=bottom, axis y line=left,
]
\addplot+[boxplot prepared={draw position=2, lower whisker=12,
  lower quartile=28.5, median=41, upper quartile=58.25, upper whisker=83,
  box extend=0.45}, color=distblue, fill=histfill!50, mark=*]
  table[row sep=\\, y index=0] {145\\ };
\addplot+[boxplot prepared={draw position=1, lower whisker=7,
  lower quartile=18, median=24, upper quartile=34, upper whisker=49,
  box extend=0.45}, color=targetgreen, fill=targetgreen!25, mark=*]
  table[row sep=\\, y index=0] {60\\ 72\\ };
\end{axis}
\end{tikzpicture}
\caption{Old and new latencies. The new median sits 17 ms lower. Is that gap
bigger than resampling noise?}
\label{fig:bootstrap-two-box}
\end{figure}

The two samples were drawn independently, so the bootstrap imitates that:
\begin{enumerate}
  \item Resample 40 values with replacement from the old group; take the median.
  \item Resample 35 values with replacement from the new group; take the median.
  \item Record the difference (old minus new). Repeat $B$ times.
\end{enumerate}

\begin{lstlisting}[language=Python]
new = np.array([32,38,27,19,24,18,33,72,24,22,43,40,34,18,31,49,14,24,
                10,15,11,28,15,38,35,29,7,23,31,34,13,24,18,20,60])
rng = np.random.default_rng(1)
m_old = np.median(rng.choice(lat, (2000, 40)), axis=1)
m_new = np.median(rng.choice(new, (2000, 35)), axis=1)
d = m_old - m_new                      # observed difference: 41 - 24 = 17
print(d.std(ddof=1), np.quantile(d, [0.025, 0.975]))   # 5.96 [ 5. 27.5]
\end{lstlisting}

\begin{figure}[H]
\centering
\begin{tikzpicture}
\begin{axis}[
  width=12cm, height=5.8cm,
  xmin=-6, xmax=38, ymin=0, ymax=340,
  axis lines=left, xlabel={bootstrap difference in medians, old $-$ new (ms)},
  ylabel={count (of 2000)}, xtick={-4,0,...,36},
  clip=false,
]
\addplot[ybar interval, fill=structurepurple!40, draw=structurepurple]
coordinates {(-4,0) (-2,2) (0,12) (2,12) (4,37) (6,76) (8,153) (10,162) (12,281)
(14,268) (16,246) (18,247) (20,144) (22,101) (24,85) (26,131) (28,30) (30,9)
(32,3) (34,1) (36,0)};
\draw[gray, thick, dashed] (axis cs:0,0) -- (axis cs:0,300);
\node[gray, anchor=south, font=\small] at (axis cs:0,300) {no difference};
\draw[densityred, very thick] (axis cs:17,0) -- (axis cs:17,310);
\node[densityred, anchor=south, font=\small] at (axis cs:17,310) {observed 17};
\draw[distblue, thick, {Stealth}-{Stealth}] (axis cs:5,220) -- (axis cs:27.5,220);
\node[distblue, anchor=west, font=\small] at (axis cs:28,220) {$[5, 27.5]$};
\end{axis}
\end{tikzpicture}
\caption{Bootstrap distribution of the difference in medians. The $95\%$
percentile interval $[5, 27.5]$ excludes 0; only 7 of the 2000 replicates
($0.35\%$) are at or below zero.}
\label{fig:bootstrap-diff-hist}
\end{figure}

Report: ``median latency dropped by 17 ms ($95\%$ CI 5 to 27.5 ms).'' The basic
interval is $[34 - 27.5,\ 34 - 5] = [6.5, 29]$; both exclude zero. The
interval is wide because the two medians wobble on their own: in these
resamples the old median has SE $5.22$ and the new one $2.91$. SEs of
independent estimates combine by squares, $\sqrt{5.22^2 + 2.91^2} = 5.98$,
matching the $5.96$ above. The bootstrap does this combination for you, for
any statistic.

\begin{trapbox}
Do not pool the two groups and resample from the pool to build a CI. Pooling
pretends both groups come from one population, which is the no-difference
hypothesis. That is a valid move for a \emph{test} (Section~\ref{sec:boot-perm})
but it gives the wrong spread for an interval around the observed difference.
\end{trapbox}

\section{An exact interval for the median from order statistics}
\label{sec:boot-exact}

For medians and other quantiles there is something better than the bootstrap:
an interval whose coverage is guaranteed for \emph{any} continuous population
and any $n$, with no resampling at all. It needs only one fact.

\begin{intuitionbox}
For a continuous population, each observation lands below the true median $m$
with probability exactly $1/2$. That is the definition of the median. So the
number of observations below $m$ behaves like the number of heads in $n$ fair
coin flips.
\end{intuitionbox}

Write $x_{(1)} \le x_{(2)} \le \dots \le x_{(n)}$ for the sorted data. These
are the \textbf{order statistics}; $x_{(j)}$ is the $j$-th smallest value. Let
$K$ be the number of observations below $m$. Then $K \sim \Bin(n, 1/2)$
(Probability From the Ground Up, Ch.~7).

Now ask when the interval $[x_{(j)}, x_{(n+1-j)}]$ misses $m$.
\begin{itemize}
  \item The whole interval sits above $m$ ($m < x_{(j)}$) exactly when fewer
  than $j$ observations are below $m$: $K \le j - 1$.
  \item The whole interval sits below $m$ ($m > x_{(n+1-j)}$) exactly when
  fewer than $j$ observations are above $m$: $n - K \le j-1$.
\end{itemize}
By the symmetry of $\Bin(n,1/2)$ both events have the same probability, and they
cannot happen together. So:

\begin{theorem}[Distribution-free CI for the median]
If $X_1,\dots,X_n$ are iid from a continuous distribution with median $m$, then
\[
  \P\!\left( X_{(j)} \le m \le X_{(n+1-j)} \right) = 1 - 2\,\P(K \le j-1),
  \qquad K \sim \Bin(n, \tfrac12).
\]
\end{theorem}

Choose the largest $j$ that keeps the coverage at least $95\%$. For $n=20$ the
table below lists the binomial probabilities.

\begin{center}
\small
\begin{tabular}{c c c c c}
\toprule
$k$ & $\P(K = k)$ & $\P(K \le k)$ & interval $[x_{(k+1)}, x_{(20-k)}]$ & coverage $1 - 2\P(K\le k)$\\
\midrule
2 & 0.0002 & 0.0002 & $[x_{(3)}, x_{(18)}]$ & 0.9996\\
3 & 0.0011 & 0.0013 & $[x_{(4)}, x_{(17)}]$ & 0.9974\\
4 & 0.0046 & 0.0059 & $[x_{(5)}, x_{(16)}]$ & 0.9882\\
5 & 0.0148 & 0.0207 & $[x_{(6)}, x_{(15)}]$ & \textbf{0.9586}\\
6 & 0.0370 & 0.0577 & $[x_{(7)}, x_{(14)}]$ & 0.8847\\
7 & 0.0739 & 0.1316 & $[x_{(8)}, x_{(13)}]$ & 0.7368\\
8 & 0.1201 & 0.2517 & $[x_{(9)}, x_{(12)}]$ & 0.4966\\
\bottomrule
\end{tabular}
\end{center}

The answer for $n=20$ is $[x_{(6)}, x_{(15)}]$, with exact coverage $95.86\%$.
Figure~\ref{fig:bootstrap-binom20} shows where the number comes from.

\begin{figure}[H]
\centering
\begin{tikzpicture}
\begin{axis}[
  width=12cm, height=5.6cm,
  xmin=-0.7, xmax=20.7, ymin=0, ymax=0.2,
  axis lines=left, xlabel={$K$ = number of the 20 observations below $m$},
  ylabel={probability}, xtick={0,2,...,20},
  yticklabel style={/pgf/number format/fixed, /pgf/number format/precision=2},
  ybar, bar width=0.6, clip=false,
]
\addplot[fill=histfill, draw=distblue] coordinates {(6,0.0370) (7,0.0739)
(8,0.1201) (9,0.1602) (10,0.1762) (11,0.1602) (12,0.1201) (13,0.0739)
(14,0.0370)};
\addplot[fill=trapred!60, draw=trapred] coordinates {(0,0.0000) (1,0.0000)
(2,0.0002) (3,0.0011) (4,0.0046) (5,0.0148) (15,0.0148) (16,0.0046)
(17,0.0011) (18,0.0002) (19,0.0000) (20,0.0000)};
\node[trapred, font=\small, align=center] at (axis cs:2.6,0.06)
  {$K \le 5$: miss\\prob.\ 0.0207};
\node[trapred, font=\small, align=center] at (axis cs:17.4,0.06)
  {$K \ge 15$: miss\\prob.\ 0.0207};
\end{axis}
\end{tikzpicture}
\caption{$\Bin(20, 1/2)$. The interval $[x_{(6)}, x_{(15)}]$ misses the median
only when 5 or fewer, or 15 or more, observations fall below it (red). Total
miss probability $2 \times 0.0207 = 0.0414$.}
\label{fig:bootstrap-binom20}
\end{figure}

\begin{walkbox}
Take the first 20 of our latencies. Sorted:
\[
  \begin{array}{cccccccccc}
  12 & 18 & 21 & 23 & 24 & \mathbf{24} & 33 & 39 & 40 & 42\\
  43 & 48 & 50 & 53 & \mathbf{63} & 64 & 68 & 68 & 70 & 79
  \end{array}
\]
The sample median is $(42+43)/2 = 42.5$. The $95.86\%$ interval is
$[x_{(6)}, x_{(15)}] = [24, 63]$ ms. No resampling, no assumption about the
shape of the latency distribution, and the coverage is exact rather than
approximate.
\end{walkbox}

\begin{figure}[H]
\centering
\begin{tikzpicture}[x=0.135cm]
\draw[-{Stealth}] (5,0) -- (88,0) node[right, font=\small] {ms};
\foreach \t in {10,20,...,80} { \draw (\t,0.1) -- (\t,-0.1)
  node[below, font=\footnotesize] {\t}; }
\foreach \v/\h [count=\i] in {12/1,18/1,21/1,23/1,24/1,24/2,33/1,39/1,40/1,42/1,
  43/2,48/1,50/1,53/1,63/1,64/2,68/1,68/2,70/1,79/1} {
  \ifnum\i<6 \fill[gray!60] (\v,0.3*\h) circle (3pt);
  \else\ifnum\i>15 \fill[gray!60] (\v,0.3*\h) circle (3pt);
  \else \fill[distblue] (\v,0.3*\h) circle (3pt); \fi\fi }
\draw[targetgreen, line width=2.5pt] (24,1.3) -- (63,1.3);
\node[targetgreen, above, font=\small] at (43.5,1.4) {$[x_{(6)}, x_{(15)}] = [24, 63]$};
\draw[densityred, thick] (42.5,-0.3) -- (42.5,1.0);
\node[densityred, anchor=north, font=\footnotesize] at (42.5,-0.45) {median 42.5};
\node[gray, font=\footnotesize] at (14,1.1) {5 below};
\node[gray, font=\footnotesize] at (73,1.1) {5 above};
\end{tikzpicture}
\caption{The order-statistic interval counts in from each end: drop 5
points on each side and keep the range of the middle 10 (blue). Stacked dots
are ties or near-ties.}
\label{fig:bootstrap-orderstat}
\end{figure}

For all 40 latencies the same recipe with $\Bin(40, 1/2)$ gives
$[x_{(14)}, x_{(27)}] = [32, 52]$ with coverage $96.15\%$. Compare the
bootstrap percentile interval $[33, 52]$: nearly the same. That is no
accident. A resample median is at most $x_{(j)}$ roughly when at least half of
the 40 draws land on the $j$ smallest data values. Each draw does so with
probability $j/40$, so that count is $\Bin(40, j/40)$. The bootstrap's lower
endpoint is where this event has probability $2.5\%$. With the normal
approximation, that is where half the draws, 20, sits 1.96 SDs above the mean
$j$:
\[
  20 \approx j + 1.96\sqrt{40 \cdot \tfrac12 \cdot \tfrac12}
  \quad\Longrightarrow\quad j \approx 20 - 6.2 = 13.8 .
\]
The exact method chose the largest $j$ with $\P(K \le j-1) \le 0.025$ for
$K \sim \Bin(40, 1/2)$: again a cut-off about $1.96 \times 3.16 = 6.2$ below
the centre 20, which the exact table turns into $j = 14$. Both methods
are counting coin flips with probability near $1/2$, so for the median they
land on nearly the same order statistics. The exact method gets there with a
guarantee; the bootstrap gets there by simulation.

The guarantee assumed a continuous population. If the population is discrete
(latencies rounded to whole ms, with ties), each observation is at or below $m$
with probability at least $1/2$ and at or above $m$ with probability at least
$1/2$. Both miss probabilities can only shrink, so the closed interval
$[x_{(j)}, x_{(n+1-j)}]$ covers at least as often as the table says.

\subsection{Any quantile}

The argument works for any quantile $q_p$ (the value with fraction $p$ of the
population below it). Now each observation falls below $q_p$ with probability
$p$, so $K \sim \Bin(n, p)$, and
\[
  \P\!\left( X_{(j)} \le q_p < X_{(k)} \right) = \P(j \le K \le k-1),
  \qquad K \sim \Bin(n,p).
\]
For a rough first guess, use the normal approximation to the binomial:
$j \approx np - 1.96\sqrt{np(1-p)}$ and $k \approx np + 1.96\sqrt{np(1-p)} + 1$,
then adjust with exact binomial probabilities.

\begin{walkbox}
p90 from $n=40$: $np = 36$, $\sqrt{np(1-p)} = 1.90$, so the guess is
$j \approx 32.3$, $k \approx 40.7$. The upper index runs past $n$. Exact search:
$[x_{(32)}, x_{(40)}] = [63, 145]$ has coverage
$\P(32 \le K \le 39) = 0.970$. The upper end is the sample maximum.
\end{walkbox}

That interval is honest about what 40 points can tell you about a p90: the
upper limit is ``somewhere up to the largest thing we saw.'' The bootstrap
percentile interval $[59.4, 83.0]$ looked tighter only because the bootstrap
cannot see beyond the data. For an upper bound on p99 the arithmetic is stark:
the sample max exceeds $q_{0.99}$ with probability $1 - 0.99^n$, which first
reaches $95\%$ at $n = 299$. With fewer than 299 requests, no order statistic
gives a distribution-free one-sided $95\%$ upper bound on the p99.

\begin{recoverybox}
Pause and think: with $n = 5$ observations, can any pair of order statistics
give a $95\%$ distribution-free interval for the median? (Answer: no. The widest
choice is $[x_{(1)}, x_{(5)}]$, which misses when all five land on one side of
$m$: probability $2 \times (1/2)^5 = 1/16$, so coverage is only $93.75\%$. With
$n = 6$ the range $[x_{(1)}, x_{(6)}]$ reaches $1 - 2/64 = 96.9\%$. Six is the
smallest sample that supports a $95\%$ median interval at all.)
\end{recoverybox}

\section{How many resamples? Monte Carlo error}
\label{sec:boot-B}

There are two sources of randomness in a bootstrap answer.
\begin{enumerate}
  \item \textbf{Sampling error}: the data are one random sample. This is what
  the interval is supposed to describe. More resamples cannot reduce it; only
  more data can.
  \item \textbf{Monte Carlo error}: we used $B$ random resamples instead of
  all $n^n$ possible ones. Re-running with a different seed gives a slightly
  different answer. This shrinks like $1/\sqrt{B}$.
\end{enumerate}

To see the second kind, repeat the whole 40-latency bootstrap 300 times with
different seeds for each $B$ and look at how much the reported SE moves
(Figure~\ref{fig:bootstrap-mcerror}).

\begin{figure}[H]
\centering
\begin{tikzpicture}
\begin{axis}[
  width=12cm, height=6cm,
  xmode=log, xmin=20, xmax=12000, ymin=4, ymax=6.5,
  axis lines=left, xlabel={number of resamples $B$ (log scale)},
  ylabel={bootstrap SE of median},
  xtick={25,100,1000,10000}, xticklabels={25,100,1000,10000},
  legend pos=north east, legend style={font=\small},
]
\addplot[name path=up, draw=none, forget plot] coordinates {(25,6.218) (50,6.058)
(100,5.785) (200,5.585) (500,5.511) (1000,5.444) (2000,5.395) (5000,5.356)
(10000,5.333)};
\addplot[name path=lo, draw=none, forget plot] coordinates {(25,4.261) (50,4.622)
(100,4.808) (200,4.927) (500,5.088) (1000,5.157) (2000,5.184) (5000,5.220)
(10000,5.239)};
\addplot[bootorange!35] fill between[of=up and lo];
\addlegendentry{middle 90\% of 300 re-runs}
\addplot[bootorange!80!black, thick, mark=*, mark size=1.5pt] coordinates {(25,5.252)
(50,5.321) (100,5.296) (200,5.282) (500,5.302) (1000,5.293) (2000,5.287)
(5000,5.286) (10000,5.288)};
\addlegendentry{average over re-runs}
\end{axis}
\end{tikzpicture}
\caption{Monte Carlo error in the bootstrap SE. With $B = 25$ one run might
report anything from 4.3 to 6.2; with $B = 1000$ the band is 5.16 to 5.44;
with $B = 10{,}000$ it is 5.24 to 5.33. The centre does not move: more
resamples sharpen the estimate of $\SE(\hat F_n, n)$ but never change it.}
\label{fig:bootstrap-mcerror}
\end{figure}

For a roughly bell-shaped bootstrap distribution, the Monte Carlo SD of the
bootstrap SE is about $\SE/\sqrt{2B}$: at $B=1000$, about $2.2\%$ of the SE.
Here we measured $0.09$ out of $5.29$, about $1.7\%$. Percentile endpoints
need more resamples because they depend on the few replicates in each tail:
at $B = 1000$ the lower $2.5\%$ endpoint had Monte Carlo SD $0.37$ ms; at
$B=10{,}000$, $0.10$ ms.

\begin{center}
\small
\begin{tabular}{l l}
\toprule
goal & resamples $B$\\
\midrule
bootstrap SE, rough & 200\\
bootstrap SE, reported to 2 digits & 1000\\
percentile or BCa $95\%$ interval & 2000--10{,}000\\
$99\%$ interval or small p-value & 10{,}000+\\
\bottomrule
\end{tabular}
\end{center}

Computers are fast and the latency median of 10{,}000 resamples takes a
fraction of a second. Pick $B$ big enough that re-running with a new seed
does not change the digits you report. Then fix and report the seed.

\section{How well do the intervals cover?}
\label{sec:boot-coverage}

A $95\%$ interval method should capture the truth in $95\%$ of repeated
samples (Chapter 7). We can check this by simulation because we know the
population that generated our latencies: log-normal with median exactly 40 ms.
For each sample size $n$, draw 4000 samples, build each interval (with
$B = 1000$ for the bootstrap ones), and count how often it contains 40.

\begin{figure}[H]
\centering
\begin{tikzpicture}
\begin{axis}[
  width=12cm, height=7cm,
  xmode=log, xmin=8, xmax=100, ymin=0.78, ymax=1.0,
  xtick={10,20,40,80}, xticklabels={10,20,40,80},
  axis lines=left, xlabel={sample size $n$ (log scale)},
  ylabel={actual coverage of nominal 95\% CI},
  yticklabel style={/pgf/number format/fixed, /pgf/number format/precision=2},
  legend style={font=\small, at={(0.5,-0.22)}, anchor=north, column sep=0.25cm},
  legend columns=3,
]
\addplot[targetgreen, thick, mark=square*] coordinates {(10,0.979) (20,0.959) (40,0.957) (80,0.966)};
\addlegendentry{exact order-statistic}
\addplot[distblue, thick, mark=*] coordinates {(10,0.938) (20,0.942) (40,0.938) (80,0.946)};
\addlegendentry{percentile}
\addplot[structurepurple, thick, dashed, mark=o] coordinates {(10,0.923) (20,0.940) (40,0.934) (80,0.943)};
\addlegendentry{BCa}
\addplot[residgray, thick, mark=triangle*] coordinates {(10,0.930) (20,0.939) (40,0.930) (80,0.932)};
\addlegendentry{$\hat\theta \pm 1.96\,\widehat{\SE}_{\text{boot}}$}
\addplot[quotientorange, thick, mark=diamond*] coordinates {(10,0.806) (20,0.827) (40,0.854) (80,0.869)};
\addlegendentry{basic}
\addplot[black, dashed, domain=8:100] {0.95};
\end{axis}
\end{tikzpicture}
\caption{Coverage of five $95\%$ intervals for the median of a log-normal
population (4000 simulated samples per point; Monte Carlo SD about $0.0035$).
The exact interval never drops below $95\%$; its zig-zag comes from the
binomial steps (Section~\ref{sec:boot-exact}). Percentile and BCa run 1 to 2
points low and creep up slowly with $n$. The basic interval under-covers badly
for the median.}
\label{fig:bootstrap-coverage}
\end{figure}

Four lessons from Figure~\ref{fig:bootstrap-coverage}:
\begin{itemize}
  \item For a median or quantile, use the exact order-statistic interval. It
  is guaranteed, and it is at least as cheap as the bootstrap.
  \item The bootstrap percentile interval is a good general default, but for
  the median at these sample sizes expect it to run 1 to 2 points short of
  nominal (about $94\%$ instead of $95\%$).
  \item BCa buys nothing for the median (the jackknife issue above). For smooth
  statistics it is usually the best bootstrap interval.
  \item The basic interval is not a safe default. Its mirror-image logic fails
  when the bootstrap distribution is lumpy and skewed.
\end{itemize}

\section{When the bootstrap fails}
\label{sec:boot-fails}

The bootstrap needs two things: $\hat F_n$ must resemble $F$ in the ways that
matter for the statistic, and the statistic must respond smoothly to small
changes in the data. Each failure below breaks one of them.

\subsection{Tiny samples}

With $n = 8$ there are only $\binom{15}{8} = 6435$ distinct resamples (as
multisets, ignoring order) and the median takes a few dozen values. $\hat F_n$
is a bag of 8 slips; it cannot represent the population's spread or tail.
Below $n \approx 15$ treat any bootstrap interval as rough, and prefer an exact
method when one exists.

\subsection{Extremes: the maximum}

Resample the 40 latencies and take the max. The largest value, 145, is in the
resample unless all 40 draws miss it:
\[
  \P(\max{}^* = 145) = 1 - (39/40)^{40} = 0.637 .
\]
Figure~\ref{fig:bootstrap-max} shows the result: almost two-thirds of the
bootstrap distribution is a single spike at the observed max.

\begin{figure}[H]
\centering
\begin{tikzpicture}
\begin{axis}[
  width=10cm, height=5.4cm,
  ybar, bar width=0.8cm, ymin=0, ymax=0.75,
  symbolic x coords={$\le 70$,79,83,145}, xtick=data,
  axis lines=left, xlabel={bootstrap maximum (ms)}, ylabel={probability},
  nodes near coords, nodes near coords style={font=\small,
    /pgf/number format/fixed, /pgf/number format/precision=3},
  enlarge x limits=0.2,
]
\addplot[fill=trapred!50, draw=trapred] coordinates {($\le 70$,0.015) (79,0.114)
(83,0.235) (145,0.637)};
\end{axis}
\end{tikzpicture}
\caption{Exact bootstrap distribution of the maximum of 40 latencies. The
replicate can never exceed the observed maximum, and $63.7\%$ of the time it
equals it. Any interval built from this distribution sits at or below 145, even
though the population maximum is larger.}
\label{fig:bootstrap-max}
\end{figure}

The problem is structural. If data come from $\Unif(0, \theta)$ and you want
$\theta$, every bootstrap replicate is at most $x_{(n)} < \theta$, so the
percentile interval has coverage exactly zero. For the max, min, and very high
quantiles, $\hat F_n$ is wrong in exactly the region the statistic looks at.

\subsection{Heavy tails}

A \textbf{heavy-tailed} distribution is one where extreme values are common
enough that the variance is huge or infinite. The Pareto family with tail
index $\alpha$ is the standard example; smaller $\alpha$ means heavier tails,
and the variance is infinite for $\alpha \le 2$. Bootstrap percentile
intervals for the \emph{mean} of $n = 50$ Pareto values (2000 simulations,
$B = 1000$):

\begin{figure}[H]
\centering
\begin{tikzpicture}
\begin{axis}[
  width=12cm, height=5.6cm,
  ybar, bar width=0.9cm, ymin=0, ymax=1.08,
  symbolic x coords={Pareto 3,Pareto 1.5,Pareto 1.1,AR(1) iid,AR(1) block},
  xtick=data, x tick label style={font=\small, align=center},
  axis lines=left, ylabel={coverage of nominal 95\%},
  nodes near coords, nodes near coords style={font=\small, anchor=north,
    yshift=-2pt, /pgf/number format/fixed, /pgf/number format/fixed zerofill,
    /pgf/number format/precision=2},
  enlarge x limits=0.13, clip=false,
]
\addplot[fill=bootorange!60, draw=bootorange!80!black] coordinates {(Pareto 3,0.889)
(Pareto 1.5,0.704) (Pareto 1.1,0.261) (AR(1) iid,0.585) (AR(1) block,0.876)};
\draw[black, dashed] (rel axis cs:0,0.8796) -- (rel axis cs:1,0.8796);
\node[anchor=south east, font=\small] at (rel axis cs:1,0.885) {target 0.95};
\end{axis}
\end{tikzpicture}
\caption{Where the percentile bootstrap breaks. Left three bars: CI for the
mean of 50 Pareto values with tail index 3, 1.5, 1.1. Right two: CI for the
mean of 200 autocorrelated values, resampling single points versus blocks of
10. Each bar is 1000--2000 simulated datasets.}
\label{fig:bootstrap-failures}
\end{figure}

With $\alpha = 1.1$ coverage collapses to $26\%$. The sample rarely contains
the giant values that dominate the population mean, so $\hat F_n$ has a much
lighter tail than $F$ and the bootstrap confidently reports a small spread.
For heavy-tailed metrics (revenue per user, file sizes, p99.9 latency) prefer
the median or a trimmed mean, or a log transform, over bootstrapping the raw
mean.

\subsection{Dependent data and the block bootstrap}

Resampling individual points assumes they were drawn independently. A latency
time series is not: a slow minute is usually followed by another slow minute.
For a series where each value is $0.7$ times the previous one plus noise (an
AR(1) process, ``autoregressive of order 1''), $n = 200$, the ordinary
bootstrap CI for the mean covered only $58.5\%$ of the time. Shuffling points
destroys the runs of high and low values, so it underestimates how much the
mean of a correlated series wanders.

The \textbf{moving block bootstrap} resamples contiguous chunks instead
(Figure~\ref{fig:bootstrap-blocks}): cut the series into all overlapping blocks
of length $\ell$, draw $n/\ell$ blocks with replacement, and glue them
together. Correlation inside each block survives. With $\ell = 10$, coverage
rose to $87.6\%$, still short of $95\%$ but far closer. Longer blocks keep more
dependence but leave fewer distinct blocks; a common rule of thumb is $\ell$
around $n^{1/3}$, tuned up for strongly correlated series.

\begin{figure}[H]
\centering
\begin{tikzpicture}[
  blk/.style={draw, minimum height=0.55cm, font=\footnotesize, inner sep=1pt},
]
\node[anchor=east, font=\small] at (-0.2,0) {series};
\foreach \c/\lab [count=\i] in {sourceblue/A,targetgreen/B,quotientorange/C,
  structurepurple/D,trapred/E,distblue/F} {
  \node[blk, minimum width=1.6cm, fill=\c!25] at (1.65*\i-0.8,0) {block \lab}; }
\node[anchor=west, font=\footnotesize] at (9.9,0) {time $\to$};
\node[anchor=east, font=\small] at (-0.2,-1.1) {resample 1};
\foreach \c/\lab [count=\i] in {quotientorange/C,sourceblue/A,quotientorange/C,
  trapred/E,targetgreen/B,structurepurple/D} {
  \node[blk, minimum width=1.6cm, fill=\c!25] at (1.65*\i-0.8,-1.1) {block \lab}; }
\node[anchor=east, font=\small] at (-0.2,-2.0) {resample 2};
\foreach \c/\lab [count=\i] in {distblue/F,distblue/F,targetgreen/B,
  structurepurple/D,sourceblue/A,trapred/E} {
  \node[blk, minimum width=1.6cm, fill=\c!25] at (1.65*\i-0.8,-2.0) {block \lab}; }
\end{tikzpicture}
\caption{Block bootstrap: the unit of resampling is a run of consecutive
observations, so short-range dependence is kept inside each block. (Shown with
non-overlapping blocks for clarity; the moving block version allows every
starting point.)}
\label{fig:bootstrap-blocks}
\end{figure}

The same logic applies to clustered data. If you have 50 users each with 20
sessions, resample \emph{users}, not sessions. The unit you resample must be
the unit that was sampled independently.

\section{Bootstrap vs permutation: a preview}
\label{sec:boot-perm}

Section~\ref{sec:boot-two-groups} asked ``how big is the drop, give or take?''
Chapter 9 asks a different question: ``could a drop this big appear if the
deploy changed nothing?'' Both are answered by resampling, but differently
(Figure~\ref{fig:bootstrap-vs-perm}).

\begin{figure}[H]
\centering
\begin{tikzpicture}[
  box/.style={draw, rounded corners, align=center, font=\small,
    text width=5.4cm, minimum height=1cm},
  arr/.style={-{Stealth[length=2.5mm]}, thick},
]
\node[box, fill=gray!10, text width=5cm] (d) at (3.3,0) {old (40) and new (35) latencies};
\node[box, fill=bootorange!15] (b) at (0,-1.7) {\textbf{Bootstrap}: resample each
  group \emph{with} replacement, separately};
\node[box, fill=structurepurple!15] (p) at (6.6,-1.7) {\textbf{Permutation}: pool all
  75, shuffle labels, split 40/35 \emph{without} replacement};
\node[box, fill=bootorange!15] (b2) at (0,-3.4) {spread of the difference\\
  $\Rightarrow$ confidence interval};
\node[box, fill=structurepurple!15] (p2) at (6.6,-3.4) {difference if labels
  meant nothing\\$\Rightarrow$ p-value};
\draw[arr] (d) -- (b); \draw[arr] (d) -- (p);
\draw[arr] (b) -- (b2); \draw[arr] (p) -- (p2);
\end{tikzpicture}
\caption{Two resampling tools with different jobs. The bootstrap keeps the
groups apart and measures uncertainty around the observed effect. The
permutation test enforces ``no difference'' by shuffling labels and asks how
unusual the observed effect is under that assumption.}
\label{fig:bootstrap-vs-perm}
\end{figure}

\begin{connectionbox}
The bootstrap shows up throughout machine learning. \emph{Bagging} (bootstrap
aggregating) trains one model per bootstrap resample and averages them;
random forests are bagged decision trees. The $\approx 37\%$ of points left out
of each resample (the $1/e$ from Section~\ref{sec:boot-resample}) form the
``out-of-bag'' set, a free validation set. Evaluation code that reports a
model's accuracy ``$\pm$'' something often bootstraps the test set: resample
test examples, recompute accuracy, take percentiles. Chapter 15 returns to
this.
\end{connectionbox}

\begin{historybox}
Bradley Efron, a statistician at Stanford, published ``Bootstrap methods:
another look at the jackknife'' in 1979. The \emph{jackknife} (Maurice
Quenouille, 1949; John Tukey, 1958) estimated bias and standard errors by
recomputing a statistic with one observation left out at a time, but it gave
nonsense for the median and offered no way to get a full sampling
distribution. Efron saw that the jackknife was a crude approximation to a
cleaner idea: treat the sample as the population and simulate. The name comes
from the phrase ``pulling yourself up by your own bootstraps,'' since the data
are used to assess themselves. The idea only became practical as computers got
cheap; it unlocked standard errors and intervals for medians, correlations,
and complex estimators that had no workable formula, and Efron's later BCa
method (1987) made the intervals accurate enough for routine use.
\end{historybox}

\section{Choosing a method}
\label{sec:boot-choose}

\begin{figure}[H]
\centering
\begin{tikzpicture}[
  q/.style={draw, rounded corners, fill=sourceblue!10, align=center,
    font=\small, text width=4.8cm, minimum height=0.9cm},
  a/.style={draw, fill=targetgreen!12, align=center, font=\small,
    text width=4.3cm, minimum height=0.9cm},
  w/.style={draw, fill=trapred!12, align=center, font=\small,
    text width=4.3cm, minimum height=0.9cm},
  arr/.style={-{Stealth[length=2mm]}, thick},
]
\node[q] (q1) at (0,0) {Observations independent?};
\node[w] (blk) at (6.8,0) {block or cluster bootstrap};
\node[q] (q2) at (0,-2.1) {Statistic is a max, min,\\or extreme quantile?};
\node[w] (ext) at (6.8,-2.1) {bootstrap fails;\\get more data or a model};
\node[q] (q3) at (0,-4.2) {Median or quantile\\of a single sample?};
\node[a] (ex) at (6.8,-4.2) {exact order-statistic interval};
\node[q] (q4) at (0,-6.3) {Has a stable SE formula\\per resample?};
\node[a] (bt) at (6.8,-6.3) {bootstrap-t};
\node[a] (bca) at (0,-8.4) {BCa (or percentile)};
\draw[arr] (q1) -- node[above, font=\footnotesize] {no} (blk);
\draw[arr] (q1) -- node[left, font=\footnotesize] {yes} (q2);
\draw[arr] (q2) -- node[above, font=\footnotesize] {yes} (ext);
\draw[arr] (q2) -- node[left, font=\footnotesize] {no} (q3);
\draw[arr] (q3) -- node[above, font=\footnotesize] {yes} (ex);
\draw[arr] (q3) -- node[left, font=\footnotesize] {no} (q4);
\draw[arr] (q4) -- node[above, font=\footnotesize] {yes} (bt);
\draw[arr] (q4) -- node[left, font=\footnotesize] {no} (bca);
\end{tikzpicture}
\caption{A decision path for a $95\%$ interval when no closed-form formula
applies. With $n$ below about 15, read any bootstrap answer as rough.}
\label{fig:bootstrap-decision}
\end{figure}

\section{Summary}
\label{sec:boot-summary}

\begin{itemize}
  \item \textbf{Plug-in principle}: the empirical distribution $\hat F_n$ (each
  data point with weight $1/n$) stands in for the unknown population.
  \item \textbf{Resample}: draw $n$ values with replacement; compute the
  statistic; repeat $B$ times. The replicates form the bootstrap distribution,
  centred near $\hat\theta$, with roughly the spread and shape of the sampling
  distribution.
  \item \textbf{Bootstrap SE} = SD of the replicates. For the mean it
  reproduces $\hat\sigma/\sqrt n$.
  \item \textbf{Intervals}: percentile $[\hat\theta^*_{(.025)}, \hat\theta^*_{(.975)}]$
  (good default); basic $[2\hat\theta - \hat\theta^*_{(.975)}, 2\hat\theta - \hat\theta^*_{(.025)}]$
  (often worse); bootstrap-t (studentized, for statistics with a stable SE);
  BCa (bias and skew corrected, best for smooth statistics).
  \item \textbf{Two groups}: resample each group separately, record the
  difference.
  \item \textbf{Exact median CI}: $[x_{(j)}, x_{(n+1-j)}]$ with coverage
  $1 - 2\P(K \le j-1)$, $K \sim \Bin(n, 1/2)$; for $n=20$ use $j=6$
  ($95.86\%$). For quantile $q_p$ use $K \sim \Bin(n, p)$.
  \item \textbf{Choose $B$} so the reported digits do not change with the seed:
  1000 for an SE, 2000--10{,}000 for intervals.
  \item \textbf{Failures}: tiny $n$, maxima and extreme quantiles, heavy tails,
  dependent data (use blocks or clusters).
  \item \textbf{Bootstrap vs permutation}: bootstrap gives intervals around an
  effect; permutation gives p-values under ``no effect'' (Chapter 9).
\end{itemize}

Chapter 9 turns the question around: instead of ``what range of values is
plausible?'' it asks ``is this one value, zero, plausible?'' and builds the
machinery of hypothesis tests.

\section{Exercises}
\label{sec:boot-exercises}

\begin{enumerate}
  \item The data are $\{3, 5, 9\}$. How many distinct resamples are there,
  counting a resample as a multiset? What is the probability that a resample's
  median is 9?
  \item For $n = 10$, what is the probability that a given observation is
  absent from a resample? Compare with the limit as $n \to \infty$.
  \item Use $\Bin(10, 1/2)$ to find the order statistics that give a
  distribution-free CI for the median with coverage at least $95\%$, and state
  the exact coverage.
  \item A statistic has $\hat\theta = 10$ and bootstrap percentiles
  $\hat\theta^*_{(0.025)} = 8$, $\hat\theta^*_{(0.975)} = 15$. Write down the
  percentile and the basic intervals. Which side of $\hat\theta$ is longer in
  each?
  \item With $B = 200$ resamples you get a bootstrap SE of 5.0. Roughly how
  much would that number change if you re-ran with a new seed? How large must
  $B$ be to make the Monte Carlo SD about $0.05$?
  \item Data are iid $\Unif(0, \theta)$, $n = 20$. What is the probability that
  a bootstrap resample's maximum equals the sample maximum? What is the coverage
  of the percentile interval for $\theta$?
  \item What is the smallest $n$ for which the sample maximum is a one-sided
  $95\%$ upper confidence bound for the p95?
  \item You have 2000 user sessions from 100 users, and want a CI for mean
  session length. What do you resample, and why?
  \item To bootstrap the correlation between height and weight from 50 people,
  a colleague resamples 50 heights and, independently, 50 weights. What will
  the bootstrap distribution look like, and what is the fix?
\end{enumerate}

\subsection*{Answers}

\begin{enumerate}
  \item $\binom{3+3-1}{3} = 10$ multisets. The median is 9 when at least two of
  three draws are 9: $3(1/3)^2(2/3) + (1/3)^3 = 7/27 \approx 0.259$.
  \item $(0.9)^{10} = 0.349$; the limit is $1/e = 0.368$.
  \item $\P(K \le 1) = 11/1024$, so $[x_{(2)}, x_{(9)}]$ has coverage
  $1 - 22/1024 = 0.979$. ($[x_{(3)}, x_{(8)}]$ gives only $0.891$.)
  \item Percentile $[8, 15]$, longer on the right (5 vs 2). Basic
  $[20 - 15, 20 - 8] = [5, 12]$, longer on the left.
  \item Monte Carlo SD $\approx 5.0/\sqrt{400} = 0.25$. Need
  $5/\sqrt{2B} = 0.05$, so $2B = 10{,}000$, $B = 5000$.
  \item $1 - 0.95^{20} = 0.642$. Coverage is $0$: every replicate is at most
  $x_{(n)} < \theta$, so the interval lies entirely below $\theta$.
  \item Need $1 - 0.95^n \ge 0.95$, i.e.\ $0.95^n \le 0.05$:
  $n \ge \ln 0.05 / \ln 0.95 = 58.4$, so $n = 59$.
  \item Resample users (all of a user's sessions together). Sessions from one
  user are correlated; resampling sessions treats 2000 dependent values as
  independent and gives an interval that is too narrow.
  \item Pairing is destroyed, so replicates are correlations of unrelated
  variables, centred near 0 with spread about $1/\sqrt{50}$. Resample whole
  $(\text{height}, \text{weight})$ pairs.
\end{enumerate}
